% VERDICT — project documentation
% Paste straight into Overleaf. Compiles with pdfLaTeX; no exotic packages.
\documentclass[11pt,a4paper]{article}

\usepackage[T1]{fontenc}
\usepackage[utf8]{inputenc}
\usepackage[margin=2.4cm]{geometry}
\usepackage{microtype}
\usepackage{booktabs}
\usepackage{longtable}
\usepackage{array}
\usepackage{enumitem}
\usepackage{xcolor}
\usepackage{tcolorbox}
\usepackage{listings}
\usepackage{amsmath}
\usepackage{graphicx}
\usepackage[hidelinks]{hyperref}

\definecolor{accent}{HTML}{D9202B}
\definecolor{ink}{HTML}{171413}
\definecolor{muted}{HTML}{6B6461}
\definecolor{paper}{HTML}{F7F5F2}

\lstset{
  basicstyle=\ttfamily\footnotesize,
  backgroundcolor=\color{paper},
  frame=single,
  rulecolor=\color{muted!40},
  breaklines=true,
  columns=fullflexible,
  showstringspaces=false,
  xleftmargin=4pt,
  framexleftmargin=4pt,
}

\newtcolorbox{keybox}[1]{
  colback=paper, colframe=accent!70, boxrule=0.6pt, arc=2pt,
  left=6pt, right=6pt, top=5pt, bottom=5pt,
  title={\bfseries #1}, coltitle=white, colbacktitle=accent
}

\newcommand{\code}[1]{\texttt{\small #1}}

\title{\textbf{VERDICT}\\[0.3em]
\large A Network Anomaly Root-Cause Assistant\\[0.2em]
\normalsize\color{muted} How the pipeline actually works, end to end}
\author{}
\date{}

\begin{document}
\maketitle
\thispagestyle{empty}

\begin{abstract}
\noindent
VERDICT takes the noisy telemetry of a microservice system during an incident ---
metrics, logs, alerts, configuration changes --- and answers one question:
\emph{which component actually broke?} It does not merely guess. It produces a
ranked list of hypotheses, each carrying the evidence behind it, an honest label
saying how strongly that evidence supports it, and a rehearsed fix. Two engines
investigate --- one arithmetic, one agentic --- and may be run simultaneously and
fused, but the arithmetic always decides. The design principle running through
every stage is that the system should be able to say ``I don't know'' and say
\emph{why}. This document explains what happens inside, in plain language, and is
candid about the parts that do not work as well as one might hope --- including
a bug that made the system announce it had found nothing on every run it ever
performed, and which it was, strictly speaking, right to announce.
\end{abstract}

\vspace{1em}

\section{The problem, in one page}

When a website breaks, the symptoms are everywhere and the cause is somewhere.
A database quietly runs out of connections; within seconds the service that
queries it starts timing out, the service above that returns errors, and the
user-facing page throws 500s. An engineer's dashboard now shows five red
services. Four of them are victims. One is the culprit. They all look equally
sick.

Worse, other things happen at the same time. Somebody pushed a config change two
minutes ago --- unrelated, but it \emph{looks} guilty because it came first.
Traffic rose 12\%. A noisy alert fired. Human root-cause analysis is largely the
work of discarding these.

VERDICT automates that discarding. Its central claim is not ``we can find the
root cause'' --- lots of systems claim that. It is:

\begin{keybox}{The actual claim}
Every conclusion is traceable to evidence, every conclusion carries a label
saying how strong that evidence is, and when the evidence is not strong enough,
the system says so instead of picking a winner anyway.
\end{keybox}

\subsection{Why this is hard}

Three reasons, each of which the pipeline attacks directly.

\begin{description}[leftmargin=1.4em, style=nextline]
  \item[Correlation is not causation.] Everything that broke, broke at roughly
  the same time. Time alone cannot separate cause from effect.
  \item[The graph matters, and it points the ``wrong'' way.] If \code{catalogue}
  calls \code{catalogue-db}, then a failure \emph{in the database} shows up
  \emph{in the catalogue}. Causes flow \emph{against} the direction of the
  arrows. Getting this backwards inverts the entire model, silently.
  \item[Absence of evidence is not evidence of absence.] If a component emits no
  telemetry, it looks perfectly healthy. The quietest thing on the dashboard may
  be the thing that is broken.
\end{description}

\section{The shape of the system}

The pipeline is a sequence of stages. Each one narrows the field, and each one
writes down \emph{why} in a permanent record called the \textbf{ledger}.

\begin{center}
\footnotesize
\begin{tabular}{@{}r@{\ \ }l@{\ \ }p{9.4cm}@{}}
\toprule
\textbf{Stage} & \textbf{Module} & \textbf{What it does} \\
\midrule
1 & \code{ingest/}    & Turn four kinds of raw telemetry into one uniform event shape. \\
2 & \code{detect/}    & Decide what counts as abnormal at all. \\
3 & \code{ledger/seed} & File every detected anomaly as a fact --- what we \emph{saw}, before any conclusion drawn from it. \\
4 & \code{localize/}  & Work out the blast radius --- which region of the graph is involved. \\
5 & \code{rank/}      & Score every plausible suspect and order them. \\
6 & \code{rank/counterfactual} & Delete a suspect and ask whether the incident still makes sense. \\
7 & \code{twin/}      & Simulate the fault and check the simulation matches reality. \\
8 & \code{rank/tiers} & Stamp each hypothesis \textsc{confirmed}, \textsc{correlated}, or \textsc{missing\_evidence}. \\
9 & \code{agents/}    & Optionally, an LLM investigates, attacks its own answer, and rehearses a fix. \\
10 & \code{rank/ensemble} & Optionally, run both engines at once and fuse their two verdicts into one. \\
11 & \code{narrate/}   & Write the report, where every claim must cite a ledger fact. \\
12 & \code{api/}      & Stream the whole thing to a browser as it happens. \\
13 & \code{chat/}     & Answer questions about a finished run, grounded in that run's ledger. \\
\bottomrule
\end{tabular}
\end{center}

\begin{keybox}{One scorer, two ways of feeding it}
The most useful thing to hold in your head: the ``agentic'' and ``deterministic''
modes are \emph{not two models}. They are the \emph{same} scorer run over
\emph{different evidence}. The agent chooses which expensive checks to buy; the
autopilot always buys the same sweep. Both end in \code{scorer.rank}. Everything
interesting about the difference between them --- including the ensemble in
\S\ref{sec:ensemble} --- follows from that one fact.
\end{keybox}

\subsection{The edge convention (the most important paragraph here)}

An edge $A \rightarrow B$ means ``$A$ calls $B$''. Therefore:

\begin{align}
  \text{symptoms that suspect } S \text{ explains} &= \operatorname{ancestors}(G, S) \cup \{S\} \\
  \text{possible causes of a sick } X &= \operatorname{descendants}(G, X) \cup \{X\}
\end{align}

Read that twice. If \code{catalogue-db} breaks, the things that \emph{show}
symptoms are its ancestors: \code{catalogue}, then \code{front-end}. The cause
is \emph{below}; the pain is \emph{above}.

This is so easy to invert --- and an inversion is so quiet, producing plausible
but exactly wrong answers --- that the repository contains a test whose only job
is to scream if anyone ever flips it (\code{tests/test\_convention.py}). It
checks the convention against \code{networkx} and against the live scorer, on the
grounds that a docstring cannot be executed.

\section{Stage by stage}

\subsection{Ingest: making four things look like one thing}

Metrics, logs, alerts and config changes arrive in different shapes from
different systems. Everything is normalised into one \code{EventEnvelope}:
who (\code{component\_id}), when (\code{ts}), what kind (\code{source}), and a
payload specific to the kind.

One rule matters more than it sounds: \textbf{component names are normalised in
exactly one place}. Real telemetry names the same service five ways
(\code{catalogue}, \code{catalogue-7d9f}, \code{catalogue.prod}, and so on). If
two stages disagree about what a component is called, the graph silently
develops holes. So one function owns it, and nothing else is allowed to clean up
a name by hand.

\subsection{Detect: what counts as broken?}

Four detectors, one per modality:

\begin{center}
\footnotesize
\begin{tabular}{@{}ll p{7.6cm}@{}}
\toprule
\textbf{Modality} & \textbf{Method} & \textbf{Idea} \\
\midrule
metric & \code{mad\_zscore} & How many robust standard deviations from this metric's own baseline? \\
metric & \code{isolation\_forest} & Does this combination of metrics look unlike normal, even if each is individually fine? \\
log    & \code{log\_freq\_spike} & Is a message template suddenly appearing far more often? \\
log    & \code{log\_rare\_template} & Did a message appear that almost never appears? \\
alert  & \code{alert\_dedup} & Collapse an alert storm into one signal. \\
config & \code{config\_risky\_flag} & Did someone change something dangerous, like a connection-pool size? \\
\bottomrule
\end{tabular}
\end{center}

A component is only called anomalous if it clears an \emph{aggregation gate}:
two independent signals, or one severe one. This is what stops a single twitchy
metric from nominating a suspect.

\subsection{Localize: the blast radius}

Given the sick components, the system takes a $k$-hop neighbourhood around them
--- the region of the graph the incident could plausibly live in. Everything
outside is ignored, which keeps the search small and the scoring fair.

\subsection{Rank: five questions about every suspect}

Every component that could be the cause gets scored on five factors, each
between 0 and 1. Let $A$ be the set of anomalous components, and let
$R(S) = \operatorname{ancestors}(G,S) \cup \{S\}$ be everything suspect $S$
could break. Then $E(S) = R(S) \cap A$ is what $S$ actually explains.

\begin{center}
\footnotesize
\begin{tabular}{@{}l l p{6.9cm}@{}}
\toprule
\textbf{Factor} & \textbf{Roughly} & \textbf{The question it asks} \\
\midrule
coverage       & $|E(S)| / |A|$ & Of everything that broke, how much does $S$ explain? \\
topo\_consistency & (see \S\ref{sec:known}) & Does the damage match the shape of the graph? \\
precedence     & $1 - \text{viol}/|E(S)|$ & Did $S$ break \emph{before} its supposed victims? An effect cannot precede its cause. \\
corroboration  & $|\text{modalities}|/4$ & Do several \emph{different kinds} of signal agree? \\
pagerank       & $\mathrm{PPR}(S)/\max \mathrm{PPR}$ & Is $S$ structurally central to the sick region? \\
\bottomrule
\end{tabular}
\end{center}

The five are combined with weights that sum to $1$, giving a score in $[0,1]$.

\begin{keybox}{The arithmetic contract}
The weighted contributions are what ship to the frontend, and they must
\emph{sum to the score} (within $10^{-6}$). This is not decoration: it means the
displayed bar chart is the score, not an illustration of it. The invariant is
enforced in code and asserted in tests, because JSON Schema cannot express it.
\end{keybox}

The weights live in one file and are tuned \emph{only} on a development split of
the data, by one script, which writes its grid search to
\code{eval/tuning\_log.json}. Held-out cases are never touched during tuning.
This is what makes the reported accuracy meaningful rather than self-congratulatory.

\subsection{Counterfactual: the idea that makes it work}

Scoring alone cannot separate a cause from a coincidence. So the system performs
an experiment:

\begin{keybox}{Remove the suspect and see if the story still holds}
Take suspect $S$ out of the graph. Ask: can the \emph{remaining} suspects still
explain the anomalies? If yes, $S$ was never necessary --- it is
\textbf{redundant}, and its score is halved. If no, $S$ is \textbf{load-bearing}
--- the incident does not make sense without it.
\end{keybox}

Concretely: ``removing \code{payment} still leaves 100\% of anomalies explained''
means \code{payment} is doing no explanatory work. ``Removing \code{catalogue}
leaves only 25\% explained'' means \code{catalogue} is carrying the incident.
The threshold is 70\%. This single mechanism is what defeats the red herring, and
it is the reason the system can \emph{clear} a suspect rather than merely
ranking it lower.

Note the careful wording. The system says \textbf{counterfactual-unchanged}. It
never says \textbf{innocent} --- innocence is ground truth, and ground truth is
reserved for offline evaluation only. Confusing ``we tested it and it did no
work'' with ``it is innocent'' would be exactly the overclaim the project exists
to avoid.

\subsection{The twin: simulate it and see}

VERDICT contains a small discrete-event simulation (SimPy) of the service graph.
It injects the hypothesised fault into the simulation, runs it, and compares the
resulting pattern of latency/error deltas against what really happened, using
cosine similarity.

\begin{center}
\footnotesize
\begin{tabular}{@{}lll@{}}
\toprule
\textbf{Similarity} & \textbf{Verdict} & \textbf{Meaning} \\
\midrule
$\geq 0.80$ & \textsc{match}    & Simulating this fault reproduces the incident. \\
$\geq 0.50$ & \textsc{partial}  & It reproduces some of it. \\
$< 0.50$    & \textsc{mismatch} & Simulating this fault produces a different incident. \\
\bottomrule
\end{tabular}
\end{center}

The twin also reports \textbf{missing evidence}: components where the simulation
predicted a symptom but reality has no instrumentation to confirm it. That is a
list of things worth wiring up before the next outage.

\subsection{Tiers: the honest label}

This is the heart of the product. Each hypothesis is stamped with one of three
labels, by exactly one piece of code, using four conditions:

\begin{center}
\footnotesize
\begin{tabular}{@{}l p{10.2cm}@{}}
\toprule
\textbf{Tier} & \textbf{Awarded when} \\
\midrule
\textsc{confirmed} & \emph{All four} hold: cited evidence resolves; a topology path
connects the suspect to every observed symptom; temporal precedence holds; and the
twin verdict is \textsc{match}. And no predicted symptom sits on an uninstrumented
component. \\
\textsc{correlated} & Otherwise, if the suspect co-occurs with the symptoms. The
label carries a reason naming exactly which condition blocked confirmation. \\
\textsc{missing\_evidence} & Otherwise --- or, short-circuiting everything above,
whenever a predicted symptom lands on a component with no telemetry. \\
\bottomrule
\end{tabular}
\end{center}

Two things deserve emphasis.

First, \textbf{the uninstrumented check runs first and overrides everything}. A
suspect with a perfect twin match still lands in \textsc{missing\_evidence} if
one of its predicted symptoms is on a blind component. The system refuses to
claim confirmation it cannot support.

Second, the tier comes with its own explanation in the backend's own words ---
for example, \emph{``\textsc{correlated}: co-occurring, but confirmation blocked
by twin partial''}. The interface displays that sentence verbatim rather than
paraphrasing it.

\subsection{Agents: bounded, optional, and measured}

Optionally, three LLM agents participate:

\begin{description}[leftmargin=1.4em, style=nextline]
  \item[Investigator] Chooses which expensive checks to buy, then files findings.
  \item[Challenger] Attacks the top hypothesis, trying to break it. An attack only
  stands if its cited event \emph{resolves} and \emph{pertains} --- same component,
  upstream of it, or inside one of its anomaly windows.
  \item[Fix-Rehearsal] Proposes remedies and rehearses each in the twin.
\end{description}

Three design rules make this safe:

\begin{enumerate}[leftmargin=1.6em]
  \item \textbf{Agents cannot mutate anything except by appending a fact to the
  ledger.} Every other tool is read-only or sandboxed.
  \item \textbf{Budgets are enforced in code, not in the prompt.} The investigator
  gets 16 tool calls, 7 ``cost points'', 180 seconds. A twin costs 2 points, a
  counterfactual 1, a rehearsal 1; everything else is free. Asking a model nicely
  to be frugal is not a budget; a counter that raises is.
  \item \textbf{If any agent fails, times out, or exhausts its budget, a
  deterministic autopilot takes over and the run still completes.} The agents are
  an enhancement, never a dependency.
\end{enumerate}

\begin{keybox}{The uncomfortable measurement}
On the synthetic suite, the agentic path scores \textbf{precision@1 $= 0.39$}
against the deterministic autopilot's \textbf{$0.52$}. On these cases the LLM
makes the answer \emph{worse}. The project measures this and reports it rather
than quietly disabling the comparison --- which is, arguably, the most useful
result in the whole evaluation.
\end{keybox}

\subsection{Narration: every claim cites a fact}

The report is generated with one tool available to it: a query against the
ledger. It cannot cite an event id, because it never sees an event --- the ledger
is its entire world.

Then the citations are validated. If a claim cites a fact that does not resolve,
\textbf{the whole line is deleted} from the report. Not flagged --- deleted. The
model gets one retry with the violation quoted back to it.

There is a real asymmetry here worth knowing: a line with \emph{no} citation
passes through untouched. The mechanism polices citations that are present and
wrong; it does not force every sentence to be cited.

\begin{keybox}{What ``the ledger is its entire world'' cost, once}
For most of this project's life the narrator opened every report with:
\code{\#\# Timeline --- No timeline facts were recorded for this run.} On every
run that had ever executed.

Detection was fine; the anomalies were on screen. But \emph{nothing ever filed
them}. \code{anomaly\_observed} had zero callers in the entire repository, as did
\code{config\_change\_observed}. Every fact in a finished ledger was a
\emph{conclusion} --- counterfactual, twin, score --- and never an
\emph{observation}. On \code{red\_herring\_config-01} the detectors find twelve
anomalies; the ledger held fourteen facts and not one of them was an observation.

So the narrator looked at its world, found no anomalies in it, and said so. It
was not hallucinating an absence. It was \emph{correctly refusing to assert what
it could not cite} --- and the citation validator would have deleted the line
anyway. Filing evidence had been quietly delegated to the agent's
\code{file\_finding}; the deterministic floor, which writes its own conclusions
directly, was never given the same job. With no key, or an agent that simply
does not bother, nobody files anything.

The fix (\code{ledger/seed.py}) transcribes every detected anomaly into a fact
\emph{before} the agent runs. The same run now narrates as an actual timeline:
the three innocent config changes first, then \code{catalogue} latency
$|z| = 1136$, then the cascade --- each line carrying a fact id that resolves.
\end{keybox}

\noindent
Two smaller consequences of that fix are worth recording. First, the chat
(\S\ref{sec:chat}) is built on the same ledger, so it inherited both the
blindness and the cure without a line of its own changing. Second, seeding
changes the ledger digest at agent start, and the transcript cache is keyed on
that digest --- so every warmed transcript became stale the moment this landed.

\section{The ensemble: both engines, then arithmetic}
\label{sec:ensemble}

The measurement above invites a lazy conclusion --- ``the LLM loses, delete it''
--- which the per-scenario breakdown refutes. Neither engine dominates:

\begin{center}
\footnotesize
\begin{tabular}{@{}lccl@{}}
\toprule
\textbf{Scenario type} & \textbf{Agentic} & \textbf{Fixed} & \textbf{Who wins} \\
\midrule
\code{confounded\_pair}     & \textbf{4/4} & 3/4 & the agent, when it can afford the sweep \\
\code{red\_herring\_config} & 2/5 & \textbf{4/5} & the autopilot, by being disciplined \\
\code{clean\_cascade}       & 2/5 & \textbf{3/5} & the autopilot \\
\code{alert\_storm}         & 0/3 & \textbf{1/3} & the autopilot \\
\bottomrule
\end{tabular}
\end{center}

\noindent
Two engines that fail on \emph{different} cases is the textbook precondition for
an ensemble. So \code{--ensemble} runs both concurrently and fuses their verdicts.

\subsection{Fuse the verdicts, not the evidence}

The obvious design is to merge both ledgers and score once. It is wrong here, and
the project's own benchmark is what says so. At spending parity the agent
\emph{loses} \code{red\_herring\_config} precisely \emph{because} it can now
afford to look at everything: the extra evidence drags innocent-but-correlated
components up the ranking.

\begin{keybox}{More evidence is not monotonically better}
Pooling the two ledgers assumes it is. Fusing the two \emph{verdicts} does not:
it keeps the autopilot's discipline at weight $(1-w)$ even when the agent's
evidence misleads. The design follows the measurement, not the intuition.
\end{keybox}

\subsection{The weight is derived, not tuned}

Let $A$ and $F$ be the two rankings. The weight on the agentic side is:

\begin{equation}
w =
\begin{cases}
  0.5 & \text{if } \operatorname{top}(A) = \operatorname{top}(F) \\[0.4em]
  \dfrac{s(A_1)}{s(A_1) + s(F_1)} & \text{otherwise}
\end{cases}
\qquad
s(h) =
\begin{cases}
3 & \text{\textsc{confirmed}} \\
2 & \text{\textsc{correlated}} \\
1 & \text{\textsc{missing\textunderscore evidence}}
\end{cases}
\end{equation}

When both engines independently land on the same suspect there is nothing to
arbitrate, so they average. When they disagree, the tie is broken by how
well-evidenced each one's \emph{own} top suspect is --- and that ladder is not a
new invention, it is the tier ordering \code{tiers.py} already publishes.

\begin{keybox}{Why there is no tuned constant}
A weight fitted to the 23 synthetic cases and then reported on those same 23
cases is overfitting in a lab coat. The dev split is \emph{empty} --- $n=1$ real
case, and 20\% of one case rounds to zero --- so a fitted weight could not be
honestly validated even in principle. A test greps the weight function for float
literals and fails if anything other than $0$, $0.5$ or $1$ appears.
\end{keybox}

Fusion blends the five sub-scores \emph{component-wise}, so the blended breakdown
still sums to the blended score --- the model validator enforces that equality to
$10^{-6}$, so a score computed separately from its parts would simply be rejected.
The fused tier is \textbf{re-derived} by \code{tiers.py} from the union of both
engines' evidence, never copied from either side. Evidence is merged, not
averaged: a twin the agent ran and the autopilot skipped is still a twin that
ran, and averaging a real verdict against ``pending'' would invent a measurement.

\subsection{The ledgers must be isolated}

Each engine's ranking is a rescore over \emph{its own} ledger. Share one and the
autopilot's five counterfactuals leak into the agent's evidence and vice versa;
the two opinions collapse into one and there is nothing left to ensemble. They
cannot even share a \code{Ledger} object: fact ids are minted from an in-memory
per-component counter, so two writers on one file mint \code{fact-catalogue-0000}
twice.

The ledger the narrator finally reads is therefore rebuilt: seeded with the
observations, then given every \emph{derived} fact both engines bought, re-filed
so each gets an unambiguous canonical id and deduplicated on
(kind, components, statement) --- both engines run the same deterministic
counterfactual, so the overlap is exact.

\subsection{One live case, and what it showed}

Measured on \code{red\_herring\_config-01} with a real key ($\$0.06$, $n=1$):

\begin{center}
\footnotesize
\begin{tabular}{@{}llc@{}}
\toprule
\textbf{Engine} & \textbf{Top-1} & \textbf{Correct?} \\
\midrule
Agentic (\code{gpt-4o}) & \code{front-end} & no --- it ranked the true fault \#4 \\
Fixed (autopilot)       & \code{catalogue} & yes \\
\textbf{Ensemble}       & \textbf{\code{catalogue}} & \textbf{yes} \\
\bottomrule
\end{tabular}
\end{center}

\noindent
The arithmetic explains itself. The agent hit \code{budget\_exhausted} before
buying the counterfactual that demotes \code{front-end}, so \code{front-end} kept
its floor score of $0.554$ where the autopilot's full sweep had cut it to $0.278$.
Averaged: $0.416$, below \code{catalogue}'s $0.480$. The ensemble rescued a wrong
agentic verdict --- not by trusting the better model, but by refusing to trust
either one completely.

\begin{keybox}{What this is not}
One case is a demonstration, not a rate. The claim ``the ensemble beats both
engines'' is \textbf{unmeasured}: it needs the full synthetic suite, which costs
roughly \$1.50 of live agent runs. Until that is done, the honest statement is
that the ensemble is \emph{motivated} by the per-scenario table and
\emph{illustrated} by one case.
\end{keybox}

\section{What the interface shows}

The frontend is a single-page React application that consumes the API and
nothing else. It renders, live over a streaming connection:

\begin{itemize}[leftmargin=1.4em]
  \item the dependency graph, lighting up as events arrive --- amber for
  anomalous, red for the leading suspect, a green ring for a component the
  counterfactual cleared, and \emph{hollow dashed} for a component with no
  telemetry;
  \item the ranked hypotheses, re-ordering themselves live as the backend
  rescores --- this is where you watch the red herring fall;
  \item the agent transcripts with the budget meter filling;
  \item the report, where every citation chip resolves to its fact on hover;
  \item the chat, where the run's own evidence answers questions about it.
\end{itemize}

\noindent
Each surface has a different authority, and conflating them is the easiest
mistake available:

\begin{center}
\footnotesize
\begin{tabular}{@{}lll@{}}
\toprule
\textbf{Tab} & \textbf{Fed by} & \textbf{Authority} \\
\midrule
Incident & \code{event\_ingested}, \code{anomaly\_detected} & observation \\
Verdict  & \code{hypothesis\_ranked}, \code{tier\_changed}  & \textbf{the verdict} \\
Agents   & \code{agent\_step}, \code{agent\_done}           & none --- it only adds evidence \\
Report   & \code{narration\_chunk}                          & language, citation-bound \\
Chat     & \code{POST /run/\{id\}/chat}                     & read path --- cannot write or decide \\
\bottomrule
\end{tabular}
\end{center}

\noindent
The agent is visible everywhere and decisive nowhere. Its findings move the
ranking only by entering the ledger and being re-scored.

\section{The chat: retrieval over the run's own evidence}
\label{sec:chat}

An engineer can ask a finished run two kinds of question --- \emph{``what should
I do about this?''} and \emph{``what's the evidence for X?''} --- and be answered
out of that run's ledger.

The corpus is the ledger, plus the ranked hypotheses and the rehearsed fixes. That
is less a design than an observation: every ledger record is \emph{already} a
short natural-language statement carrying a resolvable id --- a retrieval chunk
and its citation, written by our own code. The narration is deliberately excluded:
it is itself model output, and retrieving it would let one model cite another's
prose as evidence.

\begin{keybox}{Retrieval without embeddings}
Retrieval is TF--IDF, not a vector database. The honest reason is not the
project's no-deep-learning rule --- it is that a run's ledger is a few hundred
short statements in a vocabulary \emph{our own writers control}. Lexical search
over that is deterministic, free, instant, works with the network unplugged, and
adds no dependency. A neural retriever here would be machinery in search of a
problem.
\end{keybox}

It is a read path in every sense: no \code{file\_finding}, and it does not decide
the verdict --- the prompt says so, because a chat box next to a verdict invites
exactly the opposite assumption. Citation validation is the narrator's, imported
rather than reimplemented: an unresolvable \code{[fact-\ldots]} has its claim
deleted, and the check is applied to our \emph{own} deterministic answer too --- it
polices the claim, not the author.

\subsection{What lexical retrieval costs}

Running it broke it immediately, which is the only reason we know. \emph{``What
should I do to fix this?''} retrieved \textbf{nothing at all}: the synonym map
pointed \code{fix} at ``remediation'' while the ledger says
\code{remediation\_result}, and a word tokenizer never splits the two. The single
likeliest question an on-call engineer asks, matching zero evidence.

Three things fix it, and the third is the one that lasts:

\begin{enumerate}[leftmargin=1.6em]
  \item \textbf{Compound splitting} --- \code{remediation\_result} also emits
  \code{remediation} and \code{result}, so a plain English word reaches our
  snake\_case vocabulary without a stemmer.
  \item \textbf{Pinned context} --- the top hypothesis and the recommended fix
  enter the prompt \emph{whatever} the question was, because ``what do I do?''
  shares no words with \code{scale\_replicas} and never will.
  \item \textbf{A test that fails if a synonym points at a word no writer emits.}
  Dead weight in a synonym map is invisible until someone asks the question it was
  supposed to serve.
\end{enumerate}

Two further bugs the tests found rather than the demo: \emph{``quantum tunnelling
in the mesosphere''} retrieved evidence (stop words score above zero, so ``drop
zero-similarity chunks'' was a lie until the analyzer filtered them), and
\emph{``evidence for catalogue''} ranked a \code{payment} fact second --- the
synonym \code{evidence} $\rightarrow$ \code{observed} matched the words ``the
observed symptoms'' in unrelated prose. A synonym that collides with ordinary
English in your own writing is a false-positive generator.

\begin{keybox}{A \code{200} is not proof a model answered}
With no key, under \code{OFFLINE=1}, or once the spend cap trips, the endpoint
returns \code{mode: "deterministic"} and quotes the retrieved facts instead of
failing. Degradation is a demotion, never an error --- so a client that wants to
know whether an LLM spoke must read \code{mode}, not the status code.
\end{keybox}

\begin{keybox}{One rule governs the entire interface}
The frontend never computes a tier, a score, or a rank. It renders exactly what
the backend sent. Where the interface needs something the API does not provide,
it says so on screen rather than inventing it.
\end{keybox}

\section{How it is evaluated}

Cases are split into \textbf{dev} (for tuning) and \textbf{held-out} (for
reporting), and the tuner is forbidden from looking at held-out. Twenty-five
synthetic scenarios span seven types:
\code{clean\_cascade}, \code{red\_herring\_config}, \code{alert\_storm},
\code{confounded\_pair}, \code{missing\_telemetry}, \code{topology\_drift},
\code{ambiguous}. Real cases come from the RE2-SS dataset.

Reported metrics: \textbf{AC@1} and \textbf{AC@3} (was the true cause ranked
first / in the top three), \textbf{false-blame rate} (how often a red herring got
blamed), \textbf{median time-to-RCA}, and an efficiency table (tool calls, cost
points, wall clock).

Crucially, the benchmark also runs \textbf{ablations}: the same pipeline with the
counterfactual removed, with the twin removed, with topology removed. This is how
one distinguishes ``the system works'' from ``one part of the system works and
the rest is decoration''.

\begin{keybox}{Ground truth never leaves the evaluator}
The \code{/benchmark} endpoint strips \code{truth}, \code{rank\_of\_truth} and
\code{false\_blame} from every individual run before serving it. Aggregate rates
survive, because they are computed \emph{from} ground truth but disclose it for
no individual case. There is deliberately no ``was it right?'' badge in the
interface: the answer key is not served.
\end{keybox}

\section{Known weaknesses and open problems}
\label{sec:known}

A documentation section that lists only strengths is marketing. These are real.

\begin{description}[leftmargin=1.4em, style=nextline]

  \item[The topology factor measures the wrong thing.] \code{topo\_consistency}
  is currently $|A \cap R(S)| / |A|$ --- which is \emph{recall}, nearly a
  duplicate of \code{coverage}. It cannot penalise a ``god node'' that reaches
  everything, because reaching more only ever helps it. The intended fix is
  \emph{precision}, $|E(S)| / |R(S) \cap \text{incident}|$: of everything the
  suspect could have broken, how much actually broke? An implementation of this
  exists on a branch and is \textbf{not merged}, because swapping the factor
  changes what it means and thereby invalidates the weight vector tuned against
  the old one --- it regresses the red-herring cases until the weights are
  re-tuned on dev. This is a live piece of work, not a settled design.

  \item[The agentic path underperforms the autopilot.] See above: 0.39 versus
  0.52. The LLM's value in this system is currently explanation and challenge,
  not accuracy. The ensemble (\S\ref{sec:ensemble}) is the response to this, and
  it is not yet a rebuttal --- see the next item.

  \item[The ensemble is unmeasured.] It is \emph{motivated} by a real
  per-scenario complementarity and \emph{illustrated} by exactly one live case,
  where it rescued a wrong agentic verdict. One case is a demonstration, not a
  rate. Whether fusing beats both engines across the suite is an open question
  costing roughly \$1.50 of live runs to answer, and it is entirely possible the
  answer is no. Nothing in this document should be read as claiming otherwise.

  \item[The challenger does not run in ensemble mode.] Its penalty is applied by
  a rescore over the canonical ledger, which would re-derive every score from
  scratch and silently overwrite the fused ones --- discarding the ensemble
  immediately after computing it. Folding the penalty into the fusion is a real
  piece of work, not a line to sneak in.

  \item[Held-out is tiny.] The held-out suite has $n=1$ in places. A single case
  is an anecdote, not a rate, and the interface labels it as such rather than
  printing a confident-looking percentage.

  \item[The narration stream cannot be reassembled.] The backend chunks the
  report with \code{text.split("\textbackslash n\textbackslash n")} and emits the
  pieces --- dropping the separator. Concatenating the chunks glues every heading
  onto the previous paragraph. The frontend compensates by re-joining on a blank
  line, which is exact today, but the correct fix is one line in the backend.

  \item[Budget limits are not served.] The code that enforces budgets can report
  its own state, and nothing calls it. The interface's budget meter reconstructs
  spend by counting events and mirroring the cost table from source --- honest,
  but it will drift silently if the backend's costs change.

  \item[Several ledger fact kinds are still never written.] This used to read
  ``six of the twelve''. It is now \textbf{four}: \code{anomaly\_observed} and
  \code{config\_change\_observed} acquired a writer once it became clear that
  their absence was what made the narrator announce that nothing had been
  detected. The remaining four --- \code{anomaly\_absent}, \code{topology\_path},
  \code{temporal\_order}, \code{investigation\_note} --- have no production
  writer and exist only if an LLM files one. \code{temporal\_order} is the one
  that stings: the narrator's Timeline would read far better with explicit ``this
  config change preceded that anomaly by 45 seconds'' facts, and the scorer
  already computes precedence internally without ever writing it down. Anything
  built on the presence of these four is building on sand.

  \item[Baselines do not run.] The external comparison library is not importable
  in this environment, so the baseline row is empty. The evaluator records the
  reason rather than pretending.
\end{description}

\section{Design principles, extracted}

If the project has a thesis, it is these.

\begin{enumerate}[leftmargin=1.6em]
  \item \textbf{One writer per fact.} Tiers are assigned in exactly one file.
  Component names are normalised in exactly one function. Weights live in exactly
  one vector. Anything with two homes has two values.

  \item \textbf{Enforce in code, not in prose.} A budget in a prompt is a
  suggestion. A budget in a counter that raises is a budget.

  \item \textbf{Absence of evidence must be representable.} A predicted symptom
  on an uninstrumented component is recorded as \code{observed: null} --- neither
  true nor false --- and that alone blocks confirmation.

  \item \textbf{Measure the thing you are proud of, and publish it when it
  loses.} The ablations and the agentic-versus-fixed comparison exist to be able
  to say ``this part is not pulling its weight''.

  \item \textbf{Degrade honestly.} Agent fails $\rightarrow$ autopilot completes
  the run. Citation does not resolve $\rightarrow$ the claim is deleted. No
  telemetry $\rightarrow$ the tier drops and the report says what to instrument.
  At no point does the system fill a gap with a guess.

  \item \textbf{Let the measurement design the system.} The ensemble fuses
  verdicts rather than pooling evidence for exactly one reason: the benchmark
  showed that giving the agent \emph{more} evidence made it \emph{worse} on the
  red-herring cases. The intuition said pool; the data said don't. Where a number
  and an instinct disagree, the number is the specification.

  \item \textbf{A parameter fitted to your test set is not a parameter, it is a
  memory.} The ensemble weight is derived from the tier ladder rather than tuned,
  because with an empty dev split there is no honest way to validate a fitted
  one --- and a test greps the function to keep a future contributor from
  quietly adding one.
\end{enumerate}

\section{Running it}

\begin{lstlisting}[language=bash]
# Backend (offline: replays cached agent transcripts, makes zero API calls)
OFFLINE=1 OPENAI_API_KEY= py -m backend.main --port 8000

# Frontend against that backend
cd frontend && VITE_MOCK=0 VITE_API_BASE=http://127.0.0.1:8000 npm run dev

# Frontend alone, on recorded fixtures (no backend required)
cd frontend && VITE_MOCK=1 npm run dev

# One demo scenario, end to end (resets the run's ledger, asserts warm caches)
OFFLINE=1 make demo-2

# The ensemble: both engines concurrently, then fused
py -m backend.pipeline --case red_herring_config-01 --ensemble --top 5

# Ask a finished run a question (free; --live uses gpt-4o-mini, ~$0.0002)
py -m backend.chat.chat --run red_herring_config-01 --q "what should I do?"

# The retrieval alone, with no model involved at all
py -m backend.chat.retrieve --run red_herring_config-01 --q "why catalogue?"

# What detection filed as evidence
py -m backend.ledger.seed --case red_herring_config-01

# The gate everything must pass
make golden
\end{lstlisting}

\begin{keybox}{A note on \code{speed}}
Use \code{speed=10} for anything a human watches: roughly 34 seconds per
scenario. \code{speed=0} flushes the entire ingest burst before detection
starts, so the interface sits still and then dumps everything at once, and the
agent steps never interleave. It exists for evaluation, not for demonstration.
\end{keybox}

\section{Closing}

The interesting thing about VERDICT is not that it identifies root causes.
It is the machinery around the answer: an experiment that can clear a suspect,
a simulation that can fail to reproduce a claim, a tier that refuses to say
\textsc{confirmed}, a report that deletes its own unsupported sentences, and a
benchmark that publishes the case where the clever component loses to the simple
one.

The ensemble is the same idea aimed at the system itself. Having measured that
the clever engine loses to the simple one --- and, per scenario, that each loses
where the other wins --- the response is not to pick a favourite but to run both
and let arithmetic arbitrate, with a weight nobody was allowed to tune. On the
one live case measured so far it rescued a wrong agentic verdict; whether it does
so at a rate worth quoting is still an open question, and this document says so
rather than rounding a demonstration up into a result.

Any system can produce an answer. The work here is in producing an answer that
tells you how much to trust it --- and, when the honest answer is ``not much
yet'', in saying that out loud.

\end{document}
